% SASYR 2025 - Template based on:
% samplepaper.tex, a sample chapter demonstrating the
% LLNCS macro package for Springer Computer Science proceedings;
% Version 2.20 of 2017/10/04
%
\documentclass[runningheads]{llncs}
%
\usepackage{graphicx}
\usepackage{lipsum}
\usepackage{url}
% Used for displaying a sample figure. If possible, figure files should
% be included in EPS format.
%
% If you use the hyperref package, please uncomment the following line
% to display URLs in blue roman font according to Springer's eBook style:
% \renewcommand\UrlFont{\color{blue}\rmfamily}

\begin{document}
%
\title{Tritium Breeding Optimization In STAR Fusion Power Plant}
%
%\titlerunning{Abbreviated paper title}
% If the paper title is too long for the running head, you can set
% an abbreviated paper title here
%
\author{Rocco Lombardo\inst{1} \and
Sean Dailey\inst{1}\and
Vasil Ivakimov\inst{1}\and
Issac James\inst{1}}
%
\authorrunning{Lombardo et al.}
% First names are abbreviated in the running head.
% If there are more than two authors, 'et al.' is used.
%
\institute{Pennsylvania State University\\
\email{3/Mar/2026}\\
}
%
\maketitle              % typeset the header of the contribution
%
\begin{abstract}
With the intention for the spherical tokamak advanced reactor (STAR) to generate energy; It uses a D-T reaction to generate high-energy neutrons. 
These D-T reactions have a major advantage over other fusion reactions in the fact that they have highest cross section of occurrence, compared to other popular fusion reactions, as well as the highest power density, at $\SI{34}{\watt\per\meter\cubed\kilo\pascal\squared}$
\cite{Khodak2025}. 
However, there are multiple issues associated with the use of tritium in this reaction. The first is that tritium has a relatively short half-life of 2 years,
which makes it difficult to accumulate and store for this reaction. Second, the energy emitted in this reaction is mostly transferred to the emitted neutron,
leading to extremely high energy neutrons, 14 MeV. To resolve these issues, lithium-6 is used with a neutron source to produce both tritium and a small
number of neutrons. This occurs within the blanket, which contains the lithium-6 to breed tritium. Once the D-T reaction occurs, the kinetic energy of the 14
MeV neutrons must be converted to a more useful energy form. The current moderator of choice, within the blanket, is lead.
The primary focus of our work is optimizing the blanket's configuration to promote tritium production for this D-T reaction. The secondary focus is the
optimization of the neutron shield to protect other components and personnel.


\keywords{Tritium  \and Fusion \and Neutronics.}
\end{abstract}
%
%
%
\section{Intoduction}
\subsection{Background}
The main issue with this choice of reaction is that tritium is expensive and difficult to produce and store. However, tritium can be made by the interaction of neutrons and lithium. Since the STAR utilizes a D-T reaction to produce fast 14 MeV neutrons, with an initial injection of tritium into the plasma, the neutrons emitted by the plasma can be used to breed more tritium. Therefore a the design of a breeding blanket is important if one wants to have as much tritium as possible for the D-T reaction. 

Yet there still exists the problem of maximizing neutron-lithium interactions. Since these neutrons from the plasma are high energy, the cross section of the interaction is relatively low. To raise the cross section, the neutrons need to be moderated to a lower energy. That is why lead was chosen to be mixed in with the lithium.

Over time, the lithium in the breeding blanket will decline and need to be replaced. To do this with as little maintenance as possible, the lithium and lead within the blanket will be in a liquid state, which allows for easy control of the amount of lithium present at any point in time.
\subsection{Objectives}
This is all previous work already done by the experts at Princeton Physics Plasma Lab. Our task is to both maximize the amount of tritium created within the breeding region and develop a shield to prevent these high energy neutrons from leaving the reactor in large amounts. To do this, we have two design aspects we can change. First, the geometry of the breeding region and shield can be changed quite a bit, due to extra space around the plasma chamber. The other aspect is the material composition for both the breeding region and shield. 

\begin{table}
\caption{Table captions should be placed above the
tables.}\label{tab1}
\begin{tabular}{|l|l|l|}
\hline
Heading level &  Example & Font size and style\\
\hline
Title (centered) &  {\Large\bfseries Lecture Notes} & 14 point, bold\\
1st-level heading &  {\large\bfseries 1 Introduction} & 12 point, bold\\
2nd-level heading & {\bfseries 2.1 Printing Area} & 10 point, bold\\
3rd-level heading & {\bfseries Run-in Heading in Bold.} Text follows & 10 point, bold\\
4th-level heading & {\itshape Lowest Level Heading.} Text follows & 10 point, italic\\
\hline
\end{tabular}
\end{table}


\noindent Displayed equations are centered and set on a separate
line.
\begin{equation}
x + y = z
\end{equation}
Please try to avoid rasterized images for line-art diagrams and
schemas. Whenever possible, use vector graphics instead (see
Fig.~\ref{fig1}).

\begin{figure}
\includegraphics[width=\textwidth]{fig1.eps}
\caption{A figure caption is always placed below the illustration.
Please note that short captions are centered, while long ones are
justified by the macro package automatically.} \label{fig1}
\end{figure}

\section{Project Scope Evolution}
Describe the data gathering and process to define the project scope, describing how it evolved from the initial sponsor specification to the selected scope.

\section{Gantt Chart and Division of Work}
Document the process of activity identification given the selected scope, and give insights on how time was allocated for each activity, as well as the rationale for sequencing them. Show the initial Gantt Chart.

\section{Current Progress}
Narrative of the progress on each of the activities included in the WBS. Show if there are positive (ahead of schedule) or negative (behind schedule) deviations on the activities, and discuss the origin of such deviations.

\section{Further Plans}
Discuss any changes to the scope/WBS/timeline that will enable satisfactory completion.
%
% the environments 'definition', 'lemma', 'proposition', 'corollary',
% 'remark', and 'example' are defined in the LLNCS documentclass as well.
%

For citations of references, use:
In \cite{einstein} Einstein....

In \cite{knuthwebsite} the authors

This \cite{latexcompanion} is Latex. 
%
% ---- Bibliography ----
%
% BibTeX users should specify bibliography style 'splncs04'.
% References will then be sorted and formatted in the correct style.
%
\bibliographystyle{refs-style}
\bibliography{refs}
%
\end{document}
